\ifdefined\directlua
  \documentclass[b5paper,11pt]{ltjsbook}
\else
  \documentclass[b5paper,11pt,dvipdfmx]{jsbook}
\fi
\usepackage[margin=18mm]{geometry}
\usepackage{sectioncustomize0}
\pagestyle{plain}
\setcounter{secnumdepth}{3}

\begin{document}

\chapter{科学をひらく最初の一歩}

章見出しは、左に斜めの章番号、右に原子モデルを配したデザインです。

\section{物質の構造を見てみよう}

セクションは、楕円の番号と二本の罫線を組み合わせた大見出しです。

\subsection{原子と分子の基本}

サブセクションは、黒い番号バッジと細い横罫線で示されます。

\subsubsection{電子の配置}

サブサブセクションは、番号とタイトルを一本の太い横罫線でつなぎます。

\section{見出しが連続する場合}
\subsection{直後に subsection を置く}
\subsubsection{さらに subsubsection を続ける}

連続見出し時の前後アキが過剰にならないことを確認します。

\end{document}
